\scanpage{77}
\mspage{686}

$7^\circ)$ The groups $\underline N_\rho$, \deleted{$\unsure{\underline N'_\rho}$} etc, $\underline N_\rho^{\unsure{\natural|\sharp}}$ etc.
\margin{$\unsure{\sigma\,\mathrm D\,\sigma}$}

Let us study the subgroup
\[
(152)\qquad
\begin{gathered}
\underline N_\rho\subset\underline G_{\rho,\sigma}:\\[-2pt]
\Vert\\[-2pt]
\{\underline u\in\Gamma\underline G_{\rho,\sigma}\mid u(\rho)=\rho^{\varepsilon_3(\underline u)}\}.
\end{gathered}
\]
\margin{Provisional notation, cf. below page 697}

\deleted{We} The composite homomorphism
\diagram{\begin{tikzcd}[ampersand replacement=\&,column sep=large]
\underline N_\rho \arrow[r,hook] \arrow[rr,bend right=15,"{\underline\chi_{\deleted{\unsure{\rho}}}}"'] \& \underline G_{\rho,\sigma} \arrow[r,"\underline\chi"] \& \mathbb O_{12}^{*}
\end{tikzcd}}
is an epimorphism, \deleted{for} let us check this to reassure ourselves.

Let $\underline\mu$ be a local section of $\mathbb O_{12}^{*}$, corresponding to
$\deleted{\unsure{\underline\mu|\nu_0}}\inserted{\nu_2,\nu_3,\deleted{\unsure{\nu_0}}}\unsure{\nu_0},\mu,\nu',\nu'',\nu\deleted{\unsure{,\nu_2,\nu_3}}\ldots$
We will have
\[
\deleted{(153)}\qquad
u_{\underline\mu}(\rho)=\lambda_1^\nu\rho
=(\lambda_1^{3\nu''}\lambda_1^{-\nu_3})\rho
=\operatorname{int}(\delta^{\nu''}\tau^{\nu_3})(\rho).
\]
\deleted{let us distinguish the two cases,}

\deleted{\unsure{if} $\nu_3=0$ \unsure{$\Leftrightarrow$} i.e. $\nu=3\nu''$, we will thus have}
\[
\deleted{
\begin{gathered}
u_{\underline\mu}(\rho)=(\underbrace{\lambda_1}_{\unsure{(1-\rho_{\underline M})(\delta)\Leftarrow}})^{3\nu''}\rho
=\underbrace{\bigl((1-\rho_M)(\deleted{\unsure{\lambda_1}}\delta^{\nu''})\bigr)}_{\delta^{\nu''}\rho(\delta^{\nu''})^{-1}}\rho
=\operatorname{int}(\delta^{\nu''})(\rho)\\
\unsure{\text{i.e.}}\quad \delta=\lambda_1-\lambda_0
\end{gathered}}
\]
whence
\[
(153)\qquad
\operatorname{int}\{(\delta^{\nu''}\tau_{\unsure{1}}^{\nu_3})^{-1}\cdot u_{\underline\mu}\}(\rho)=\rho
\]
i.e.
\[
(154)\qquad
\left\{
\begin{aligned}
\underbrace{(\delta^{\nu''}\tau_{\unsure{1}}^{\nu_3})^{-1}u_{\underline\mu}}_{u}
&=\tau_{\unsure{1}}^{\nu_3}\delta^{-\nu''}u_{\underline\mu}\in\Gamma\underline N_\rho\\
\underline\chi(u)&=\deleted{\unsure{\mu}}\underline\chi\bigl((\delta^{\nu''}\tau_1^{\nu_3})^{-1}\bigr)\cdot\underbrace{\underline\chi(u_{\underline\mu})}_{\underline\mu}\\
&\hspace{-1em}\underline\chi(\tau_1^{\nu_3})=\underbrace{(-1)^{\nu_3}}_{\varepsilon_3}
\end{aligned}
\right.
\]
\margin{\uline{N.B.} We denote by $\underline\chi$ the canonical homomorphism
\diagram{\begin{tikzcd}[ampersand replacement=\&]
\underline G_{\rho,\sigma}\arrow[r] \& \mathbb O_{12}^{*}
\end{tikzcd}}}

Thus setting
\[
v=(\omega\tau')^{\nu_3}u
=\bigl(\underbrace{\omega\deleted{\unsure{\tau}}\tau'\tau}_{\sigma}\bigr)^{\nu_3}\delta^{-\nu''}u_{\underline\mu}\qquad\deleted{\text{i.e.}}
\]
\deleted{we will have}
\inserted{we will have}
\[
(154)\qquad
\underline\chi(v)=\underline\mu,\quad\text{with }v=\sigma^{\nu_3}\delta^{-\nu''}u_{\underline\mu}\in\Gamma\underline N_\rho
\]
\[
\left(\delta=\lambda_1\lambda_0^{-1},\quad
\begin{gathered}\nu_3=\nu_3(\underline\mu),\\\nu''=\nu''(\underline\mu)\end{gathered}\right).
\]
\margin{76}

\scanpage{78}

We will thus have an exact sequence
\[
(155)
\]
\diagram{\begin{tikzcd}[ampersand replacement=\&,column sep=large]
1\arrow[r,dash] \& \underline{\unsure{Z|N}}_\rho^{\deleted{\unsure{0}}}\arrow[r,dash] \& \underline N_\rho\arrow[r,dash] \& \mathbb O_{12}^{*}\arrow[r,dash]\arrow[d,equal] \& 1\\
\& \& \& \gamma_{\rho,\sigma}
\end{tikzcd}}
and
\[
\underline{\unsure{Z|N}}_\rho^{\deleted{\unsure{0}}}
=\underline{\operatorname{Centr}}_{\underline{\unsure{\mathcal G|\widetilde G}}^{+}}(\rho).
\]

We have, if
\[
(156)\qquad\left\{
\begin{aligned}
H^+_{S_0}&\overset{\mathrm{dfn}}{=}
\underline{\unsure{\mathcal G|\widetilde G}}^{+}/
\underbrace{\deleted{\unsure{\underline Z}}\underline{\unsure{\mathcal G|\widetilde G}}^{0}}_{\substack{\text{neutral}\\\text{component,}\\\simeq\underline M}}
\simeq(H^+)_{S_0}\\
\text{i.e.}\quad H^+&=\unsure{\mathcal G|\widetilde G}^{+}/M,
\end{aligned}\right.
\]
a canonical homomorphism
\[
(157)
\]
\diagram{\begin{tikzcd}[ampersand replacement=\&,column sep=large]
\underline Z_\rho\arrow[r] \& \underline{\operatorname{Cent}}_{H^+_{S_0}}(\rho)\simeq\deleted{\unsure{\underline Z}}\underline L_\rho\quad(\simeq(\mathbb Z/6\mathbb Z)_{S_0})
\end{tikzcd}}
\deleted{which is \unsure{surjective}\inserted{an epimorphism} (the kernel is equal}
which is an epimorphism (since there is even a section, coming from the inclusion $\underline L_\rho\subset\underline Z_\rho$) and whose kernel is
\[
(158)\qquad
\underline Z_\rho^0=\underline Z_\rho\cap\underline M
=\{x\in\Gamma\underline M\mid\rho_{\underline M}\underline x=x\quad\text{i.e.}\quad(\unsure{1-\rho_{\underline M}|1+\rho_{\underline M}})x\deleted{\unsure{\rho}}=0\}.
\]
This kernel is zero over $\underbrace{\operatorname{Spec}\mathbb Z}_{S_0}\setminus\{3\}$,
but of dimension $1$ in dimension $3$. Yet we do feel that the ``good'' definition of $\underline Z_\rho^{\unsure{0}}$ \deleted{is} must \deleted{be by}\inserted{correspond to} an equation
\[
\deleted{\rho_{\underline M}\underline x}\qquad
(159)\qquad \unsure{\ell|\varepsilon}(1-\rho_{\underline M})x=0\qquad\uuline{\text{in }\underline M_0},
\]
\margin{77}

\scanpage{79}
\mspage{687}

which in fact amounts \inserted{simply,} to replacing $\underline N_\rho$ defined in (152), by the schematic closure of its restriction over $S_0-\{3\}$; We thus find a better subgroup, which from now on I denote by $\underline N_\rho$, fitting into the exact sequence
\[
(159)
\]
\diagram{\begin{tikzcd}[ampersand replacement=\&,column sep=large]
1\arrow[r] \& \underline L_\rho\arrow[r,phantom,"{\unsure{\hookrightarrow|\leftarrow}}" description] \& \underline N_\rho\arrow[r,"\underline\chi_\rho"] \& \underline{\mathbb O}_{12}^{*}\arrow[r] \& 1
\end{tikzcd}}

Let us set
\[
(160)\qquad\left\{
\begin{aligned}
\underline G_{\rho,\sigma}^{\unsure{\natural|\sharp}}
&=\underbrace{\underline{\mathbb O}_{12}^{*}}_{\gamma_{\rho,\sigma}}\cdot
\underbrace{\underline{\unsure{\mathcal G|\widetilde G}}^{\unsure{\natural|\sharp}}}_{\underline M}
\subset\underline{\mathbb O}_{12}^{*}\cdot\underline{\unsure{\mathcal G|\widetilde G}}^{+}
=\underline G_{\rho,\sigma},\\
\underline N_\rho^{\unsure{\natural|\sharp}}
&=\underline N_\rho\cap\underline G_{\rho,\sigma}^{\unsure{\natural|\sharp}}
\end{aligned}\right.
\]
we then find
\[
(161)\qquad
\underline N_\rho^{\unsure{\natural|\sharp}}\cap\underline L_\rho
=\underline M\cap\underline L_\rho=\{1\}
\]
thus
\[
(161)
\]
\diagram{\begin{tikzcd}[ampersand replacement=\&,column sep=large]
\underline N_\rho^{\unsure{\natural|\sharp}}\arrow[r,hook] \& \underbrace{\underline{\mathbb O}_{12}^{*}}_{\gamma_{\rho,\sigma}}
\end{tikzcd}}
I say that the image is
\[
(162)\qquad
\underline\gamma'_{\rho,\sigma}
=\underline{\mathbb O}_{12}^{\prime *}
\overset{\mathrm{dfn}}{=}
\left\{\underline\mu\in\Gamma\underline{\mathbb O}_{12}^{*}\;\middle|\;
\begin{gathered}\varepsilon_3(\underline\mu)=1\\\text{i.e. }\nu_3(\underline\mu)=0\end{gathered}\right\}.
\]
\margin{thus we \unsure{have} $\underline N_\rho^{\unsure{\natural|\sharp}}\subset\underline Z_\rho$, it is better to denote it by $\underline Z_\rho^{\unsure{\natural|\sharp}}$.}
That the image contains $\underline{\mathbb O}_{12}^{\prime *}$ follows from (154), let us show that it is contained.

Let
\[
u=f^{-1}u_{\underline\mu},\qquad
f\in\Gamma\underline{\unsure{\mathcal G|\widetilde G}}^{\unsure{\natural|\sharp}}\inserted{=\Gamma\underline M},\quad
\underline\mu\in\Gamma\underline{\mathbb O}_{12}^{*}
\]
be a \inserted{local} section of
$\underline G_{\rho,\sigma}^{\unsure{\natural|\sharp}}
=\mathbb O_{12}^{*}\underline{\unsure{\mathcal G|\widetilde G}}^{\unsure{\natural|\sharp}}$,
let us write that it is in $\underline N_\rho$, whence
\margin{78}

\scanpage{80}

\[
\operatorname{int}(f^{-1})\underbrace{u_{\underline\mu}(\rho)}_{\lambda_1^\nu\rho}
=\rho^{\varepsilon_3}\qquad\text{i.e.}\qquad
\lambda_1^\nu\rho=\operatorname{int}(f)\rho^{\varepsilon_3}
\]
reducing modulo $\underline M$ we find
\[
\rho=\deleted{\unsure{\operatorname{int}f}}\rho^{\varepsilon_3}\qquad\text{in }H^+
\]
which implies $\varepsilon_3=1$, OK. Thus we have
\[
(163)
\]
\diagram{\begin{tikzcd}[ampersand replacement=\&,column sep=large]
\deleted{\underline Z_\rho^{\unsure{\natural|\sharp}},}\underbrace{\underline N_{\rho,\deleted{\unsure{\sigma}}}^{\unsure{\natural|\sharp}}}_{=\underline Z_\rho^{\unsure{\natural|\sharp}}}\arrow[r,"\sim"] \& \underline\gamma'_{\rho,\sigma}=\underline{\mathbb O}_{12}^{*\prime}
\end{tikzcd}}
--- thus $\inserted{\underline Z_\rho^{\unsure{\natural|\sharp}}=}\underline N_{\rho,\deleted{\unsure{\sigma}}}^{\unsure{\natural|\sharp}}$ is a \uline{partial section}, over $\underline{\mathbb O}_{12}^{\unsure{\prime*\prime|*\prime}}$, of the extension $N_{\rho,\sigma}$ of $\underline{\mathbb O}_{12}^{*}$ by $L_\rho$.

\deleted{\unsure{Let us set} then} This partial section does
\[
\deleted{(164)\qquad N_{\rho,\sigma}^{?}=\{1,\tau'\}.}
\]
not contain $\tau'$, which is a section of $\underline N_{\rho\deleted{\unsure{,\sigma}}}$. Does $\tau'$ normalize $\underline{\unsure{Z|N}}_\rho^{\unsure{\natural|\sharp}}$ --- since it normalizes $\underline N_\rho$, it would suffice for it to normalize
$\deleted{\unsure{\underline{\mathcal G}^{\natural}}}\inserted{\underline G_{\rho,\sigma}^{\unsure{\natural|\sharp}}}$,
\deleted{But \inserted{which} is obvious since $\underline{\unsure{\mathcal G|\widetilde G}}^{\unsure{\natural|\sharp}}=\underline{\unsure{\mathcal G|\widetilde G}}^0$. Thus $\tau'=\tau\sigma$, where $\tau$}

Since $\tau'=\tau\sigma$, \deleted{it amounts} and
$\tau\in\Gamma\underline{\unsure{G|\mathcal G}}_{\rho,\sigma}^{\unsure{\natural|\sharp}}$
normalizes $\underline{\unsure{G|\mathcal G}}_{\rho,\sigma}^{\unsure{\natural|\sharp}}$,
it amounts to the same thing to say that $\sigma$ normalizes
$\underline{\unsure{G|\mathcal G}}_{\rho,\sigma}^{\unsure{\natural|\sharp}}
=\underline{\mathbb O}_{12}^{*}\cdot\underline{\unsure{\mathcal G|\widetilde G}}^{\unsure{\natural|\sharp}}_{\deleted{\unsure{\rho,\sigma}}}$.
Since $\sigma$ normalizes $\underline{\unsure{\mathcal G|\widetilde G}}^{\unsure{\natural|\sharp}}=\underline{\unsure{\mathcal G|\widetilde G}}^0$ which is normal, it suffices to look at the action of $\sigma$ on $\underline{\mathbb O}_{12}^{*}\subset G_{\rho,\sigma}$, \deleted{\unsure{defined|defined}} i.e. on the $u_{\underline\mu}$ of (151). The question is whether the
\[
\sigma u_{\underline\mu}\sigma^{-1}u_{\underline\mu}^{-1}\in\Gamma\underline{\unsure{\mathcal G|\widetilde G}}^{\unsure{\natural|\sharp}},
\]
now
\margin{79}